\section{Related Work and Research Gaps}
\label{section:background}

\subsection{From Scaling to Small Models}
\label{subsec:background_evolution}
\label{subsec:background_techniques}
\label{subsec:background_families}

The Transformer \citep{vaswani2017attention} and the pre-train-then-adapt paradigm \citep{devlin2019bert,brown2020gpt3} led to a scaling race that the Chinchilla study \citep{hoffmann2022chinchilla} partly redirected: a 70B model trained on 1.4~trillion tokens outperformed the 280B Gopher, implying that most large models were undertrained and that smaller, longer-trained models could match them. Subsequent families pushed this further at small scale through curated and synthetic data \citep{gunasekar2023phi1,abdin2024phi3}, distillation from larger teachers \citep{gemma2024,team2025gemma3}, and pruning plus distillation from a larger sibling, as for Llama~3.2's 1B and 3B edge models \citep{meta2024llama32}. Three bodies of technique underlie the current generation: distillation and synthetic-data training \citep{hinton2015distilling,gunasekar2023phi1}; post-training compression and parameter-efficient adaptation \citep{frantar2023gptq,lin2024awq,hu2022lora,dettmers2023qlora}; and on-device runtimes \citep{gerganov2023llamacpp,meta2024executorch,mlc2024}. Sections~\ref{sec:compression} and~\ref{sec:deployment} review them, and \suppsec~\ref{supp:families} and~\ref{supp:background} summarize the principal families and the longer background.

\subsection{LLM--SLM Collaboration and Agentic Systems}
\label{subsec:background_collab}

A line of work studies how small and large models cooperate \citep{chen2024role_small,chen2025collab_survey}. Cascades and routers are the most relevant to agents: FrugalGPT \citep{chen2023frugalgpt} learns a cascade over commercial APIs, Hybrid~LLM \citep{ding2024hybridllm} routes by predicted difficulty, AutoMix \citep{aggarwal2024automix} escalates on self-verification, RouteLLM \citep{ong2024routellm} learns from preference data, RouterBench \citep{hu2024routerbench} benchmarks routers on 405K inference outcomes, and Division-of-Thoughts \citep{division_of_thoughts_2025} routes sub-questions between an on-device and a cloud model. Speculative decoding \citep{leviathan2023speculative,chen2023specsampling} uses a small model to accelerate a large one without delegating answers. On the agentic side, function-calling data pipelines report specialized sub-10B models among the leading entries of the Berkeley Function-Calling Leaderboard \citep{liu2024apigen,liu2025toolace,zhang2024xlam}, and agent evaluation has begun to report cost alongside accuracy \citep{kapoor2024agentsmatter,kapoor2025hal,erol2025costofpass}.

\subsection{Gaps in Existing Literature}
\label{subsec:background_gaps}

General \gls{slm} surveys \citep{lu2024slmsurvey,vannguyen2024slm_survey,wang2024comprehensive} provide technique taxonomies, application coverage, and in the case of \citet{lu2024slmsurvey} a measurement study of 70 models; \citet{gupta2025slmsurvey}, updated in 2026, survey about 160 papers on 1--8B models. None separates active from total parameters or foregrounds agentic deployment. \citet{srivastava2025thinkslm} benchmark 72 small models on 17 reasoning benchmarks with three repeats, including models derived by quantization, pruning, and distillation, and find that reasoning tracks training method and data more than scale and that quantization preserves reasoning while pruning disrupts it; it is the closest thing in the literature to the controlled evaluation whose absence Section~\ref{sec:reported} documents, but covers reasoning only. \citet{feng2025efficient_reasoning_survey} survey efficient reasoning, and \citet{chen2024mas_survey} survey multi-agent systems by application.

Two works state the agentic case directly. \citet{belcak2025slm_agentic} argue it as a position, with a conversion algorithm and estimates of how much agent traffic small models could handle, and \citet{tian2026edgeagents} argue that personal agents must move to the edge. The closest prior work is the survey of \citet{sharma2025slm_agentic}. It reviews 1--12B (occasionally about 20B) models for agentic workloads, proposes an \gls{slm}-default, \gls{llm}-fallback architecture with uncertainty-aware routing, retries, and verifier repair (their Section~VIII), proposes cost per successful task as the deciding metric, and names where large models still win: open-domain synthesis, complex multi-hop reasoning, safety-critical judgment, and long-context work (their Section~XI). We share that design and metric. It does not separate active from total parameters, grade its evidence or report contrary results, examine whether the vendor evaluation record supports comparison, or state a decision rule under which its position would be refuted, and its estimate that cost per successful task falls 10--30$\times$ (their Section~X) is not a reported measurement. Those are the points on which this paper differs: the C1/C2 definition, the verified inventory, the quantified audit of the evaluation record, graded evidence including contrary results, a break-even condition that includes retries, checking, and escalation together with a correctness requirement on the check, and a falsification design with comparators, a decision rule, and a sample size. Table~\ref{tab:survey_compare} summarizes the positioning.

\begin{table}[t]
\centering
\caption{Coverage of this survey relative to prior surveys and position papers (Y~full, P~partial, --~absent), as judged by the present authors from each work's text (grade~\evD). \emph{Active/total}: separates active from total parameters. \emph{Evaluation critique}: analysis of whether the reporting record supports cross-model comparison. \emph{Graded evidence}: an explicit strength-of-evidence grade per claim. \emph{Agentic focus}: small models inside agents as a primary topic. \emph{Decision rule}: a stated condition under which the position is refuted.}
\label{tab:survey_compare}
\footnotesize
\adjustbox{max width=\linewidth}{%
\begin{tabular}{lccccccc}
\toprule
\textbf{Work} & \textbf{Taxonomy} & \textbf{Compression} & \textbf{Active/} & \textbf{Evaluation} & \textbf{Graded} & \textbf{Agentic} & \textbf{Decision} \\
 & & & \textbf{total} & \textbf{critique} & \textbf{evidence} & \textbf{focus} & \textbf{rule} \\
\midrule
\citet{lu2024slmsurvey}            & Y & P & -- & P & -- & -- & -- \\
\citet{vannguyen2024slm_survey}    & Y & Y & -- & -- & -- & -- & -- \\
\citet{wang2024comprehensive}      & Y & Y & -- & -- & -- & P & -- \\
\citet{gupta2025slmsurvey}         & Y & P & -- & -- & -- & -- & -- \\
\citet{srivastava2025thinkslm}     & P & P & -- & Y & -- & -- & -- \\
\citet{belcak2025slm_agentic}      & -- & -- & -- & P & -- & Y & -- \\
\citet{tian2026edgeagents}         & -- & -- & -- & -- & -- & Y & P \\
\citet{sharma2025slm_agentic}      & P & P & -- & -- & -- & Y & -- \\
\textbf{This survey}               & Y & Y & Y & Y & Y & Y & Y \\
\bottomrule
\end{tabular}%
}
\end{table}
